\section{Problem and Parallelism Boundary}
\subsection{Layer-Aware Construction}
Let $G=(V,E)$ be the directed routing-resource graph (RRG), with $N=|V|$ nodes and $M=|E|$ edges. Nodes carry resource type, layer, geometry and an indexed delay/congestion cost $w(v)=(d_v,q_v)$. The graph contains planar and cross-layer channels as well as logic, memory and DSP connectivity~\cite{third_dimension,multidie}. PathFinder-style negotiation resolves contention~\cite{pathfinder}; VPR consumes lookahead costs during routing~\cite{vtr8,vtr9}. We accelerate construction of the sampled wire table~\cite{diversity}, not negotiated routing itself.

A key $g=(\ell_s,\sigma,\tau)$ identifies each group in $\mathcal G$ by source layer, segment type and channel type. Native enumeration yields $J$ jobs indexed by $a$, each retaining group $g(a)$, ordinal, source RR node $s_a=s_{g,j}$ and recording context. One search covers all target layers. Its map $\phi_a(v)$ assigns an eligible IPIN to a source-relative cell $u=(\ell_t,\Delta x,\Delta y)$ or returns $\bot$ when geometry excludes it. A grid of width $X$, height $Y$ and $L$ layers has $C=LXY$ result cells per sample.

An eligible recorded pin contributes one path's accumulated delay $d$ and indexed resource base cost $q$. A cell $R_a[u]$ remains missing if no eligible visit records it before search termination. Native insertion $\oplus$ fills a missing cell or replaces its pair on strictly smaller delay. For each group,
\begin{equation}
 A_g[u]=\mathop{\bigoplus}_{a:g(a)=g}^{\text{native order}}R_a[u],\qquad
 F=\mathcal I\bigl(\mathcal C(\{A_g\})\bigr).
\end{equation}
$\mathcal C$ and $\mathcal I$ denote native group writes and interpolation in their prescribed interleaving and missing-cell scan order. The table is indexed by $(\ell_s,\ell_t,\tau,\sigma,\Delta x,\Delta y)$. Source/OPIN preprocessing and derived minima remain in the outer computation.

\subsection{Necessary Dependencies and Removable Barriers}
\begin{figure}[t]
\centering
{\footnotesize\begin{tikzpicture}[x=1mm,y=1mm,font=\footnotesize,
  bar/.style={draw,fill=blue!9},wait/.style={draw,dashed,fill=gray!8}]
\foreach \yy/\txt in {0/Stock, -7/Group $p_1$, -13/Group $p_2$, -20/STRICT $p_1$, -26/STRICT $p_2$}
 \node[anchor=west] at (0,\yy) {\txt};
\foreach \start/\stop/\yy/\job in {19/49/0/$A_1$,49/59/0/$A_2$,59/71/0/$B_1$,19/49/-7/$A_1$,49/61/-7/$B_1$,19/29/-13/$A_2$,19/49/-20/$A_1$,19/29/-26/$A_2$,29/41/-26/$B_1$}{
 \draw[bar] (\start,\yy-2.4) rectangle (\stop,\yy+2.4);
 \node at ({(\start+\stop)/2},\yy) {\job};}
\draw[wait] (29,-15.4) rectangle (49,-10.6);
\node at (39,-13) {wait};
\draw[dotted,thick] (49,-4)--(49,-16);
\node[anchor=west] at (52,-13) {group join};
\node[anchor=west] at (0,-34) {ALT GPU};
% Illustrative sample-local keys in two successive GPU batch rounds.
\node[draw,fill=orange!10] at (32,-34) {$A_1:b,\ B_1:c$};
\node[draw,fill=orange!10] at (65,-34) {$A_1:b,\ B_1:c'$};
\draw[->,>=Stealth] (46,-34)--(51,-34);
\end{tikzpicture}
}
\caption{Illustrative schedules for three samples. With Group, a worker waits for $A_1$ before starting $B_1$. With STRICT, $A_2$ and $B_1$ finish before $A_1$; their results wait for native-order consumption. In successive ALT GPU rounds $A_1$ can reopen bucket $b$ while $B_1$ advances from $c$ to $c\prime$.}\label{fig:contracts}
\end{figure}
For point-like locations on one row, starts at $x=0$ and $5$ reaching a pin at $x=6$ record displacements 6 and 1. A multi-source search retaining only the cheaper path loses one contribution. Within a cell, equal delays preserve the earlier sample's entire pair. For native-order contributions $(5,9),(5,3),(6,1)$, the winner is $(5,9)$. Consuming the second first gives $(5,3)$, and separate field minima give $(5,1)$, a pair absent from all three paths.

In the locked Stock path, \texttt{get\_routing\_cost\_map} iterates over sampled nodes and calls a Dijkstra flood for each~\cite{lockedvtr}. Each search can instead read shared costs and maintain private labels. It can start before the preceding group joins (Fig.~\ref{fig:contracts}). Changing priority extraction inside a sample can alter its representative and requires a separate search rule.

\subsection{Lookahead Construction in FPGA Routers}
Sharma and Hauck~\cite{adaptive_astar} reduce table storage by clustering wires and searching from each cluster as a single source. Petelin and Betz~\cite{diversity} sample individual wires by type and orientation. For each displacement they record the smallest observed delay and its associated path cost. AIR uses map-based estimates to guide adaptive incremental routing~\cite{air}. The associated path cost in our baseline follows the sampled-wire representation of Petelin and Betz. For 2.5D interposer fabrics, Poolad et al.~\cite{multidie} separate the intra-die estimate from the cost of crossing each intervening die boundary.

These choices determine which costs a router can look up. Building a layer-aware map also requires scheduling the sampled searches. Singer et al.~\cite{multiqueue} evaluate parallel connection searches and suggest using parallel Dijkstra searches to sample more starting wires. We hold the sampled starts fixed. STRICT overlaps native searches while preserving their ordered contributions. ALT processes updates within each search in parallel and checks how the resulting cost pairs affect routed circuits.
